\documentclass[11pt]{article}
\pdfinfoomitdate=1
\pdftrailerid{}
\pdfsuppressptexinfo=15
\usepackage[a4paper,margin=24mm]{geometry}
\usepackage{amsmath,amssymb}
\usepackage[T1]{fontenc}
\usepackage{xcolor}
\usepackage{fancyhdr}
\usepackage{listings}
\definecolor{cobalt}{HTML}{2457D6}
\definecolor{graphite}{HTML}{17202A}
\definecolor{muted}{HTML}{667085}
\pagestyle{fancy}
\fancyhf{}
\lfoot{\scriptsize\color{muted}SYNTHETIC DEMO DOCUMENT --- raw student input}
\rfoot{\scriptsize\color{muted}Page \thepage\ of 2}
\renewcommand{\headrulewidth}{0pt}
\setlength{\parindent}{0pt}
\setlength{\parskip}{8pt}
\newcommand{\answer}[1]{\vspace{8pt}{\small\bfseries\color{cobalt}#1}\par\vspace{3pt}\hrule\vspace{10pt}}
\begin{document}
{\LARGE\bfseries\color{graphite} Alex Chen --- Submission}\par
{\small\color{muted}Student ID: DEMO-001 --- LaTeX response --- synthetic student}\par
\vspace{8pt}\hrule\vspace{14pt}

\answer{Q1 --- Calculus}
Let $u=x^2$, so $\mathrm{d}u=2x\,\mathrm{d}x$. Therefore
\[
I=\frac12\int_0^1 e^u\,\mathrm{d}u
 =\frac12\left[e^u\right]_0^1
 =\frac{e-1}{2}.
\]

\answer{Q2 --- Mechanics}
The normal force is $N=mg\cos 30^\circ$. Along the slope,
\[
ma=mg\sin 30^\circ-\mu_kN,
\qquad
a=9.81(\sin30^\circ-0.20\cos30^\circ)
  =3.20\,\mathrm{m\,s^{-2}}.
\]
Starting from rest, $v^2=2as$, so
\[
v=\sqrt{2(3.20)(3)}=4.38\,\mathrm{m\,s^{-1}}.
\]

\newpage
{\LARGE\bfseries\color{graphite} Alex Chen --- Submission}\par
{\small\color{muted}Student ID: DEMO-001 --- LaTeX response --- synthetic student}\par
\vspace{8pt}\hrule\vspace{14pt}

\answer{Q3 --- Linear algebra}
If $Ax=0$, then $A^{\mathsf T}Ax=0$. Conversely, if $A^{\mathsf T}Ax=0$, left-multiplying by $x^{\mathsf T}$ gives
\[
x^{\mathsf T}A^{\mathsf T}Ax=\lVert Ax\rVert^2=0,
\]
hence $Ax=0$. The kernels are equal, so rank--nullity gives equal ranks.

\answer{Q4 --- Programming}
\begin{lstlisting}[language=Python,basicstyle=\ttfamily\small,frame=single]
import math

def stable_softmax(xs):
    if not xs:
        return []
    exps = [math.exp(x) for x in xs]
    total = sum(exps)
    return [value / total for value in exps]
\end{lstlisting}
\end{document}
